\section{Syntax}
\label{sec:expr-proc}

\begin{figure}
  \setlength{\abovedisplayskip}{0pt}
  \setlength{\belowdisplayskip}{0pt}
  \begin{align*}
    \shortintertext{\textbf{Types}}
    \Mob \grmeq
              & \MobStationary \grmor \MobMobile
    \tag{Mobilities}                                                                           \\
    \Dir \grmeq
              & \DirLeft \grmor \DirRight \grmor \DirUnord \grmor \DirNone
    \tag{Directions}                                                                           \\
    \Eff \grmeq
              & \EffPure \grmor \EffImpure
    \tag{Effects with $\EffPure < \EffImpure$}                                                 \\
    \T \grmeq & \BUnit
    \grmor \TArr{\Mob}{\Dir,\Eff}\T\T
    \grmor \T \TPair \T
    \grmor \TVariant{\ell : \T[\ell]}
    \grmor \S
    \tag{Types}                                                                                \\
    \S \grmeq
              & \SSkip \grmor \SSeq\S\S
                 \grmor \SSend\T \grmor \SRecv\T
                \grmor \SChoice{\ell:\S[\ell]}{\ell\in L}
                \grmor \SBranch{\ell:\S[\ell]}{\ell\in L}
    \\ \grmor & \STerm \grmor \SWait \grmor \SDrop \grmor \SAcq \grmor
                              \SRec\SVar\S \grmor \SVar
    \tag{Session types}
    \\
    \Ctx \grmeq
              & \CtxEmpty \grmor \CtxVar \EVar\T \grmor \Ctx\CtxSeq\Ctx \grmor \Ctx\CtxPar\Ctx
                \tag{Typing environments}
    \\[\medskipamount]
    \shortintertext{\textbf{Expressions}}
    \C \grmeq
              & \CUnit
    \grmor \CNew
    \grmor \CFork
    \grmor \CTerm
    \grmor \CWait
    \grmor \CSend
    \grmor \CRecv
    \grmor \CDiscard
    \\\grmor &
      \CSplit \grmor \CSSplit
      \grmor \CDrop \grmor \CAcq
      \grmor \CSelect k
      \grmor \CBranch
    \tag{Constants}                                                                            \\
    \E \grmeq
              & \C \grmor \EVar \grmor \EAbs\EVar\E \grmor \EAppD[\Dir]{\E}{\E}
    \grmor \ESeq\E\E
    \grmor \EPair\E\E
    \grmor \ELetPair{\EVar}{\EVar}{\E}{\E}                                                     \\
    \grmor    & \EInj i \E
    \grmor \ECase \E {\ECaseBranch* i \EVar \E}
    \grmor \EMu[\EVar] \E
    \tag{Expressions}
    \\[\medskipamount]
    \shortintertext{\textbf{Processes}}
    \BindGroup \grmeq
              & (\EVar \grmor \BSep)^*
    \tag{Binder groups}
    \\
    \Bind \grmeq
              & \EVar^*
    \tag{Variable sequences}
    \\
    \Proc \grmeq
              & \PExp\E
    \grmor \PPar\Proc\Proc
    \grmor \PNew{\BindGroup}{\BindGroup}{\Proc}
    \tag{Processes}
  \end{align*}
  \caption{Syntax of types, expressions, processes}
  \label{fig:syntax}
\end{figure}

\Cref{fig:syntax} collects the source-language syntax used throughout
this paper. The figure is organized into three groups: types, expressions, and
processes. We write $\ELet\EVar{\E[1]}\E[2]$
as syntactic sugar for $\EApp{(\EAbs\EVar{\E[2]})}{\E[1]}$, where the
function is linear (unordered or left-ordered) and immobile.

The types group introduces three orthogonal classifiers. Mobilities
($\MobStationary$, $\MobMobile$) indicate if values are stationary or 
can be moved across process boundaries. Directions ($\DirLeft$,
$\DirRight$, $\DirUnord$, $\DirNone$) govern evaluation order and context
composition in function and pair constructs. An effect is either pure
$\EffPure$ or impure $\EffImpure$; effects are lattice-ordered by
$\EffPure < \EffImpure$. Types 
$\T$ combine ordinary functional forms (unit, arrow, pair, variant) with
session types $\S$. Session types are either $\SSkip$ or use
sequential composition $\SSeq\S\S$ as their central
connective. They include communication, branching/selection,
closing/waiting, drop/acquire (new to this work), and recursion.

Typing environments $\Ctx$ are tree-shaped and combine variable bindings using
ordered ($\Ctx\CtxSeq\Ctx$) and unordered ($\Ctx\CtxPar\Ctx$) composition. This
structure is inspired by 
\citet{DBLP:journals/bsl/OHearnP99}, but with different semantics.

The expression group separates constants $\C$ from compound expressions $\E$.
Constants cover unit, channel creation, process creation, communication primitives,
split operations, and synchronization primitives. They are better understood
together with their type in the next section.

Expressions add variables, abstraction and application, sequencing,
products, sums, and recursion. Following our directional presentation,
application carries an explicit direction annotation. We distinguish left
($\DirLeft$) and right ($\DirRight$) application syntactically. The
direction determines both the evaluation order and context position in the body
of the function. Throughout this paper, omitted
direction annotations or plain juxtaposition are understood as standard, left-to-right
CBV-application.

A binder group $\BindGroup$ is a list of channel names $\EVar$ and
separators $\BSep$. The metavariable $\Bind$ ranges over separator-free
variable sequences. We write $\BEmp$ for the empty list and
$\BWithSep{\Bind}{\BindGroup}$ for the concatenation of $\Bind$, $\BSep$, and
$\BindGroup$. Concatenation of $\BindGroup$s and $\Bind$s is written
by juxtaposition with unit $\BEmp$. We only use groups where each separator is followed by
at least one variable, but the syntax does not
enforce this invariant. Processes $\Proc$ only arise at run time. They
include a single expression, parallel composition of processes, and
restriction with two binder groups that scope over the body.


%%% Local Variables:
%%% mode: latex
%%% TeX-master: "../popl2027.tex"
%%% End:
